\documentclass{article}
\usepackage[scaled]{helvet}
\renewcommand*\familydefault{\sfdefault} 
\usepackage[margin=1in]{geometry} % set suitable margins
\usepackage{parskip}
\usepackage{graphicx}

\usepackage{amsmath}
\numberwithin{equation}{section}

\title{Heat Transfer Assignment \#5}
\author{Rodrigo Rojas}
\date{4-19-2024}

\begin{document}
\maketitle

\section{\#1}
\subsection*{Problem Statement}

Consider a copper sphere with a diameter of $2 \, \text{cm}$ and a thermocouple located at its center. The sphere is suddenly immersed in liquid nitrogen contained in a Dewar flask. The temperature response of the sphere is monitored using a digital data acquisition system that records the temperature every $0.05 \, \text{s}$. At the moment when the temperature $T = 92.5 \, \text{K}$, the maximum rate of temperature change $\frac{dT}{dt}$ is observed to be $19.8 \, \text{K/s}$. The saturation temperature of nitrogen at $1 \, \text{atm}$ pressure is $77.4 \, \text{K}$. Given the properties of copper at $92.5 \, \text{K}$ are: density $\rho = 8930 \, \text{kg/m}^3$, specific heat capacity $c = 235 \, \text{J/(kg}\cdot\text{K)}$, and thermal conductivity $k = 450 \, \text{W/(m}\cdot\text{K)}$, perform the following tasks:

\begin{enumerate}
  \item Using the lumped thermal capacity model, determine the corresponding heat transfer coefficient.
  \item Check the Biot number to ensure that the model is valid.
\end{enumerate}

\subsection*{Solution}

\subsubsection*{Given Data}
\begin{itemize}
    \item Diameter of the copper sphere, $d = 2 \, \text{cm} = 0.02 \, \text{m}$
    \item Maximum rate of temperature change, $\frac{dT}{dt} = 19.8 \, \text{K/s}$
    \item Temperature of the sphere when maximum rate of temperature change occurs, $T = 92.5 \, \text{K}$
    \item Saturation temperature of nitrogen at $1 \, \text{atm}$ pressure, $T_{\infty} = 77.4 \, \text{K}$
    \item Density of copper, $\rho = 8930 \, \text{kg/m}^3$
    \item Specific heat capacity of copper, $c = 235 \, \text{J/(kg} \cdot \text{K)}$
    \item Thermal conductivity of copper at $92.5 \, \text{K}$, $k = 450 \, \text{W/(m} \cdot \text{K)}$
\end{itemize}

\subsubsection*{Part 1: Evaluation of the Heat Transfer Coefficient}
This section uses the lumped capacitance model, which assumes that the temperature distribution inside a solid is nearly uniform, thus allowing us to consider the temperature throughout the solid as constant. We can describe the heat flux, $Q$, through the following relationship:
\begin{equation}
    Q = mc\frac{dT}{dt} = hA(T - T_{\infty})
\end{equation}
Here, $m$ represents the mass of the copper sphere, $c$ denotes the specific heat capacity, $\frac{dT}{dt}$ refers to the rate of temperature change, $h$ is the heat transfer coefficient, $A$ is the surface area, $T$ is the temperature of the sphere, and $T_{\infty}$ is the ambient temperature.

The volume and consequently the mass of the sphere are calculated using the sphere's density:
\begin{align}
    V &= \frac{4}{3}\pi (0.01 \, \text{m})^3 \\
    m &= \rho V = 8930 \, \text{kg/m}^3 \times \frac{4}{3}\pi (0.01 \, \text{m})^3
\end{align}
The sphere's surface area is:
\begin{equation}
    A = 4\pi (0.01 \, \text{m})^2
\end{equation}
Inserting these values into the heat flux equation to solve for $h$, we obtain:
\begin{equation}
    h = \frac{mc\frac{dT}{dt}}{A(T - T_{\infty})} = \frac{8930 \, \text{kg/m}^3 \times \frac{4}{3}\pi (0.01 \, \text{m})^3 \times 235 \, \text{J/(kg} \cdot \text{K)} \times 19.8 \, \text{K/s}}{4\pi (0.01 \, \text{m})^2 \times (92.5 \, \text{K} - 77.4 \, \text{K})}
\end{equation}
Using these computations, we determine the heat transfer coefficient to be approximately $9172.47 \, \text{W/(m}^2\text{K)}$.

\subsubsection*{Part 2: Confirmation of the Model Validity via Biot Number}
The Biot number, $Bi$, is a dimensionless quantity that evaluates the ratio of internal conductive resistance to the convective heat transfer resistance at the boundary. It is expressed as:
\begin{equation}
    Bi = \frac{hL_c}{k}
\end{equation}
where $L_c = \frac{V}{A}$ is the characteristic length, and $k$ represents the thermal conductivity. For the sphere, $L_c$ is determined by dividing the volume by the surface area:
\begin{equation}
    Bi = \frac{9172.47 \, \text{W/(m}^2\text{K)} \times \frac{\frac{4}{3}\pi (0.01 \, \text{m})^3}{4\pi (0.01 \, \text{m})^2}}{450 \, \text{W/(m} \cdot \text{K)}}
\end{equation}
This results in a Biot number of approximately $0.0679$. With $Bi << 0.1$, the lumped capacitance model is validated, confirming a minimal internal temperature gradient within the sphere.


\newpage
\section{\#2}
\subsection*{Problem Statement}

Consider a system where electronic components are mounted onto a finned aluminum base plate, designed to facilitate efficient heat dissipation. These components are actively cooled by a stream of air produced by a fan. Key thermal properties of the system are as follows:
\begin{itemize}
    \item The combined mass and specific heat capacity of the base plate and components give a total thermal capacity of \(5000 \, \text{J/K}\).
    \item The product of the effective heat transfer coefficient and the surface area of the base plate is \(10 \, \text{W/K}\).
\end{itemize}

Initially, both the base plate and the cooling air are at a temperature of \(295 \, \text{K}\). At this point, the electronic components start dissipating power at a rate of \(300 \, \text{W}\).
\begin{enumerate}
    \item Calculate the temperature of the base plate after \(10\) minutes.
    \item Determine the steady-state temperature of the components.
\end{enumerate}


\subsection*{Solution}

\subsubsection*{Given Data}
\begin{itemize}
    \item The sum of the products of mass and specific heat of the base plate and components, \(mc = 5000 \, \text{J/K}\)
    \item The product of heat transfer coefficient and surface area, \(hA = 10 \, \text{W/K}\)
    \item Initial temperature of the plate and cooling air, \(T_{\infty} = 295 \, \text{K}\)
    \item Power applied to the system, \(Q = 300 \, \text{W}\)
\end{itemize}

\subsubsection*{Part A: Transient Temperature Response}
The transient temperature response of the plate, \(T(t)\), can be analyzed using the first-order thermal response model. This approach is common in thermal systems where the heat capacity and heat transfer characteristics dominate the system's behavior. The temperature at any time \(t\) is given by:
\begin{equation}
    T(t) = T_{\infty} + (T_{ss} - T_{\infty})(1 - e^{-\frac{t}{\tau}})
\end{equation}
where \(T_{ss}\) is the steady-state temperature and \(\tau\) is the time constant, which represents the system's response speed to thermal changes. The time constant \(\tau\) is defined as the ratio of the system's thermal mass to its heat transfer capability:
\begin{equation}
    \tau = \frac{mc}{hA} = \frac{5000 \, \text{J/K}}{10 \, \text{W/K}} = 500 \, \text{s}
\end{equation}
Given \(t = 10\) minutes (\(600\) seconds), substituting these values into the temperature equation gives:
\begin{equation}
    T(600) = 295 \, \text{K} + (325 \, \text{K} - 295 \, \text{K})(1 - e^{-\frac{600}{500}}) \approx 315.96 \, \text{K}
\end{equation}
This calculation estimates the temperature of the plate after 10 minutes of operation under the given conditions.

\subsubsection*{Part B: Steady-State Temperature Calculation}
In the steady state, the system reaches a thermal equilibrium where the heat generated by the electronic components equals the heat dissipated to the environment. This condition can be mathematically represented by the equation:
\begin{equation}
    Q = hA(T_{ss} - T_{\infty})
\end{equation}
where \(Q\) is the power dissipation, \(hA\) is the product of the heat transfer coefficient and the surface area, and \(T_{\infty}\) is the ambient temperature. Solving for \(T_{ss}\) yields:
\begin{equation}
    T_{ss} = \frac{Q}{hA} + T_{\infty} = \frac{300 \, \text{W}}{10 \, \text{W/K}} + 295 \, \text{K} = 325 \, \text{K}
\end{equation}
This calculation determines the steady-state temperature of the components, indicating how hot the components will get during continuous operation in this setup.


\newpage
\section{\#3}
Consider a 10 m length of 2 cm internal diameter (I.D.) tube. The tube wall is maintained at a temperature of 320 K, and water enters the tube at 300 K, with an atmospheric pressure of 1 atm and a flow velocity of 3 m/s. Determine the average heat transfer coefficient and the pressure gradient inside the tube under these conditions.

\subsection*{Solution}
% Introduction to the concept of bulk mean temperature
The bulk mean temperature \( T_b \) is a concept used to approximate the average temperature within a fluid flowing inside a tube. It is calculated as the average of the water temperature entering the tube (inlet temperature) and the temperature of the tube wall. This provides a simple way to represent the average temperature of the fluid for use in further calculations.

% Equation for calculating bulk mean temperature
\begin{equation}
    T_b = \frac{T_{\text{inlet}} + T_{\text{wall}}}{2} = \frac{320 + 300}{2} = 310\, \text{K}
\end{equation}
This equation calculates \( T_b \) as the average of the inlet temperature (320 K) and the wall temperature (300 K), resulting in 310 K.

% Why thermophysical properties are evaluated at T_b
It's important to evaluate the thermophysical properties of water at this temperature because these properties (like viscosity and thermal conductivity) can change with temperature. Using the average temperature ensures that the calculations are more accurate and reflective of the conditions inside the tube.

% Definitions and significance of various thermophysical properties
\begin{equation}
\left\{
\begin{aligned}
    L &= 10\, \text{m} && \text{(Tube length: the distance water travels inside the tube)}, \\
    D &= 2\, \text{cm} && \text{(Internal diameter of the tube: the width across the tube)}, \\
    V_0 &= 3\, \text{m/s} && \text{(Velocity of water entering the tube: speed at which water enters)}, \\
    \text{Pr} &= 4.6 && \text{(Prandtl number: a measure that relates the fluid's ability to transport momentum vs heat)}, \\
    k &= 0.628\, \text{W/mK} && \text{(Themal conductivity: how well the fluid conducts heat)}, \\
    \rho &= 993\, \text{kg/m}^3 && \text{(Density: mass per unit volume of the fluid)}, \\
    \mu &= 6.95 \times 10^{-4}\, \text{Pa}\cdot\text{s} && \text{(Dynamic viscosity: resistance to flow)}, \\
    c_p &= 4174\, \text{J/(kgK)} && \text{(Specific heat capacity: heat capacity per unit mass)}
\end{aligned}
\right.
\end{equation}


% Explanation of Reynolds number and its importance
The Reynolds number \( Re_D \) is a dimensionless quantity used in fluid mechanics to predict flow regimes in different fluid flow situations. A Reynolds number greater than 2300 generally indicates that the flow is turbulent, which affects how heat and momentum are transferred in the tube.

% Calculation of Reynolds number
\begin{equation}
    Re_D = \frac{\rho V_0 D}{\mu} = \frac{993 \times 3 \times 0.02}{6.95 \times 10^{-4}} = 85726.6 \quad (\text{indicating turbulent flow})
\end{equation}
This calculation shows the flow is turbulent because the Reynolds number is significantly higher than 2300.

% Description of empirical correlations for turbulent flow
Empirical correlations are formulas derived from experimental data. They are used here to estimate the Nusselt number (Nu) and friction factor (f), which are crucial for understanding heat transfer and fluid resistance in turbulent conditions.

% Equations for friction factor and Nusselt number
\begin{equation}
    f_b = \left(0.79 \ln(\text{Re}_D) - 1.64\right)^{-2} = \left(0.79 \ln(85726.6) - 1.64\right)^{-2} = 0.0186
\end{equation}
\begin{equation}
    Nu_b = \frac{\frac{f}{8}(\text{Re}_D - 1000)\text{Pr}}{1 + 12.7 \left(\frac{f}{8}\right)^{0.5}(\text{Pr}^{2/3} - 1)} = \frac{\frac{0.0186}{8}(85726 - 1000)4.6}{1 + 12.7 \left(\frac{0.0186}{8}\right)^{0.5}(4.6^{2/3} - 1)} = 435.25
\end{equation}

% Viscosity correction due to temperature variations
When the fluid's temperature changes, so does its viscosity. To adjust for this, we use a viscosity correction factor.

% Equation for viscosity correction factor
\begin{equation}
    \frac{\mu_s}{\mu_b} = \text{(viscosity correction factor)} = \frac{\mu_{\text{at surface temperature}}}{\mu_{\text{at bulk mean temperature}}} = \frac{5.84 \times 10^{-4}}{6.95 \times 10^{-4}} = 0.84
\end{equation}

% Adjusting heat transfer coefficient and pressure gradient for temperature effects
\begin{equation}
    f = f_b \left(\frac{\mu_s}{\mu_b}\right)^{0.25} = 0.0186 \times 0.84^{0.25} = 0.0178
\end{equation}
\begin{equation}
    Nu = Nu_b \left(\frac{\mu_s}{\mu_b}\right)^{-0.11} = 435.25 \times 0.84^{-0.11} = 443.7
\end{equation}

% Final calculations for heat transfer coefficient and pressure gradient
\begin{equation}
    h = \frac{Nu \times k}{D} = \frac{443.7 \times 0.628}{0.02} = 13932\, \frac{\text{W}}{\text{m}^2 \text{K}}
\end{equation}
\begin{equation}
    \frac{\Delta P}{L} = \frac{1}{2} f \rho V_0^2 = \frac{1}{2} \times 0.0178 \times 993 \times 3^2 = 3977.2\, \text{Pa/m}
\end{equation}




\newpage
\section{\#4} 
A very thin stainless steel shim stock is bonded to a 10 cm-square Teflon plate to serve as a heater element. Water at 290 K flows as a laminar boundary along the plate. The voltage drop and current flow in the heater are measured to obtain the power dissipated; also, from the known dependence of electrical resistance on temperature, the average plate temperature can be deduced. In a particular test, the power is 1025 W, and the average plate temperature is 311.5 K. Estimate the water velocity, and the temperature distribution along the plate.


\subsection*{Solution}
First, we start by defining the dimensions and thermal properties of the system:

\begin{equation}
\left\{
\begin{aligned}
    L &= 10 \text{ cm} = 0.1 \text{ m} && \text{(length of the plate)}, \\
    T_{\infty} &= 290 \text{ K} && \text{(temperature of the water flowing over the plate)}, \\
    T_{s} &= 311.5 \text{ K} && \text{(average surface temperature of the plate)}, \\
    T_{r} &= \frac{290 + 311.5}{2} = 300.8 \text{ K} && \text{(used for some thermal calculations)}, \\
    Pr &= 5.78 && \text{(Prandtl number, a dimensionless number relevant in fluid flow)}, \\
    k &= 0.6125 \text{ W/m}\cdot\text{K} && \text{(thermal conductivity of the fluid)}, \\
    \nu &= 8.54 \times 10^{-7} \text{ m}^2/\text{s} && \text{(kinematic viscosity of the fluid)}
\end{aligned}
\right.
\end{equation}


\subsubsection*{Assumptions}
\noindent
We assume a laminar flow without significant conduction along the direction of the plate or through its backside, as it is insulated. This assumption allows us to focus solely on the heat transfer from the fluid to the plate's surface.

\subsubsection*{Heat Transfer Calculations}
The heat transfer from the fluid to the plate can be modeled using the Nusselt number ($Nu$), Reynolds number ($Re$), and Prandtl number ($Pr$) as follows:

\begin{equation}
    Nu_L = 0.68 Re_L^{\frac{1}{2}} Pr^{\frac{1}{3}} = \frac{h_c \times L}{k} \quad \text{(defining the Nusselt number using the heat transfer coefficient, \( h_c \))}
\end{equation}

\noindent
Substituting values to find specific heat transfer characteristics:

\begin{equation}
    Nu_L = \frac{h_c \times L}{k} = \frac{\frac{q_s}{T_s - T_{\infty}} \times L}{k} = \frac{\frac{Q}{A} \times L}{T_s - T_{\infty}} \times \frac{L}{k} = \frac{\left( \frac{1025}{(0.1 \times 0.1)} \right) \times 0.1}{(311.5 - 290)} \times \frac{0.1}{0.615} = 778.4
\end{equation}

\subsubsection*{Flow Characteristics}
Using the above calculations, we find the Reynolds number indicating the nature of the flow:

\begin{equation}
    Nu_L = 0.68 Re_L^{\frac{1}{2}} Pr^{\frac{1}{3}} = 778.4 \quad \text{(calculation from earlier, used to find Reynolds number)}
\end{equation}
\begin{equation}
    Re_L = \left( \frac{778.4}{0.68 \times Pr^{\frac{1}{3}}} \right)^{2} = 406807 = 4.06 \times 10^{5} \quad \text{(Reynolds number indicating laminar flow)}
\end{equation}
\begin{equation}
    Re_L = \frac{V \times L}{\nu} = 4.06 \times 10^{5} \times 8.54 \times 10^{-7}
\end{equation}
\begin{equation}
    V = \frac{3.4742 \, \text{m/s}}{0.1} = 34.742 \, \text{m/s} \quad \text{(water velocity calculated from Reynolds number)}
\end{equation}

\subsubsection*{Temperature Distribution}
Finally, we calculate the temperature distribution along the plate using the heat flux and the previously determined flow characteristics:

\begin{equation}
    Nu_x = \frac{h_x \times x}{k} = \frac{q'' \times x}{\Delta T_x} = \frac{q'' \times x}{k \Delta T_x} = \frac{q'' \times x}{k(T_x - T_{\infty})}
\end{equation}
\begin{equation}
    \frac{1}{Re_x^{\frac{1}{2}}} = \frac{V \times x}{\nu}
\end{equation}
\begin{equation}
    Nu_x = 0.453 Re_x^{\frac{1}{2}} Pr^{\frac{1}{3}} \Rightarrow T_x - T_{\infty} = \frac{q'' \times x^{\frac{1}{2}}}{0.453 k \times V^{\frac{1}{2}} \times Pr^{-\frac{1}{3}}}
\end{equation}
\begin{equation}
    T_x = 102.058 \, x^{\frac{1}{2}} + 290 \quad \text{(expression for temperature distribution along the plate)}
\end{equation}
Below is the graphical representation of the findings:
\begin{center}
    \includegraphics[width=0.7\textwidth]{image.png}
\end{center}
\newpage
\section{\#5}
SAE 50 oil flows in a 60 m-long, 25 mm-ID tube at 0.5 kg/s. The oil enters the tube at 300 K, and the walls of the tube are maintained at 370 K. The objective is to determine the exit bulk temperature of the oil.
\subsection*{Solution}
\subsubsection*{Given Data}
The following parameters are given, which describe the physical setup and initial conditions of the problem:

\begin{equation}
\left\{
\begin{aligned}
    \dot{m} &= 0.5 \, \text{kg/s} && \text{(Mass flow rate of the oil)}, \\
    D &= 25 \, \text{mm} = 0.025 \, \text{m} && \text{(Internal diameter of the tube)}, \\
    L &= 60 \, \text{m} && \text{(Length of the tube)}, \\
    T_{\text{in}} &= 300 \, \text{K} && \text{(Inlet temperature of the oil)}, \\
    T_{s} &= 370 \, \text{K} && \text{(Temperature of the tube walls)}
\end{aligned}
\right.
\end{equation}


To tackle this problem, we start by guessing the oil's outlet temperature to be 320 K. We'll check if this guess was correct later on. We calculate the average temperature during the flow:
\begin{equation}
    \frac{T_{\text{in}} + T_{\text{out}}}{2} = \frac{300 + 320}{2} = 310 \, \text{K}
\end{equation}

Next, we list the characteristics of SAE 50 oil at this average temperature:
\begin{equation}
\left\{
\begin{aligned}
    Pr &= 3400 && \text{(Prandtl number)} \\
    C_{p} &= 1950 \, \text{J/kg} \cdot \text{K} && \text{(Specific heat capacity)} \\
    \rho &= 877 \, \text{kg/m}^3 && \text{(Density)} \\
    k &= 0.1435 \, \text{W/m} \cdot \text{K} && \text{(Thermal conductivity)} \\
    \mu &= 0.25 \, \text{kg/m} \cdot \text{s} && \text{(Dynamic viscosity)}
\end{aligned}
\right.
\end{equation}

The Reynolds number helps us understand the type of flow, whether it's smooth (laminar) or chaotic (turbulent). We calculate it as follows:
\begin{equation}
    \text{Re}_{D} = \frac{4 \dot{m}}{\pi D \mu} = \frac{4 \times 0.5}{\pi \times 0.025 \times 0.25} = 101.86
\end{equation}
A Reynolds number under 2300 means the flow is smooth (laminar).

For smooth flow in a tube, we use a formula to find the Nusselt number, which measures the effectiveness of heat transfer:
\begin{equation}
    Nu_{D} = 3.66 + \frac{0.065 \frac{D}{L} Re_{D} Pr}{1 + 0.04 \left(\frac{D}{L} Re_{D} Pr\right)^{\frac{1}{3}}} = 3.66 + \frac{0.065 \times \frac{0.025}{60} \times 101.86 \times 3400}{1 + 0.04 \left(\frac{0.025}{60} \times 101.86 \times 3400\right)^{\frac{1}{3}}} = 8.125
\end{equation}

Since the oil's viscosity changes with temperature, we adjust the Nusselt number accordingly:
\begin{equation}
    Nu_{D} = Nu_{D} \left(\frac{\mu_{s}}{\mu_{b}}\right)^{-0.11} = 8.125 \left(\frac{181}{2500}\right)^{-0.11} = 10.846
\end{equation}

From this, we calculate the convective heat transfer coefficient, which tells us how effectively the tube transfers heat to the oil:
\begin{equation}
    h_{c} = \frac{k}{D} Nu_{D} = \frac{0.1435}{0.025} \times 10.846 = 62.257 \, \text{W/(m}^2 \text{K)}
\end{equation}

Using an energy balance, we find the oil's outlet temperature:
\begin{equation}
    T_{\text{out}} = T_{s} - (T_{s} - T_{\text{in}})e
\end{equation}

After simplifying, we determine the final outlet temperature to be:
\begin{equation}
    T_{\text{out}} = 318.190 \, \text{K}
\end{equation}

Our calculated temperature is close to our initial guess, confirming our initial assumptions and showing that the properties didn't significantly change, so no reevaluation is needed.


\newpage
\section{\#6}
Calculate the average heat transfer coefficient for a hailstone with a diameter of 6 mm falling at its terminal velocity through air at a temperature of 287 K and a pressure of 1 atm.

\subsection*{Solution}
\subsubsection*{Given Data:}
\begin{equation}
\left\{
\begin{aligned}
    d &= 6 \, \text{mm} && \text{(Diameter of the hailstone)} \\
    T &= 287 \, \text{K} && \text{(Temperature of the air)} \\
    P &= 1 \, \text{atm} && \text{(Pressure of the air)} \\
    \rho_s &= 917 \, \text{kg/m}^3 && \text{(Density of the hailstone, assumed to be constant)} \\
    \rho &= 1.2 \, \text{kg/m}^3 && \text{(Density of air, standard density at sea level)}
\end{aligned}
\right.
\end{equation}


\subsubsection*{Calculating Terminal Velocity:}
To calculate the terminal velocity (\(V_t\)) of a hailstone, we need to consider the speed it reaches when the air resistance stops it from accelerating further. For a spherical object like a hailstone, this velocity is determined by balancing the gravitational pull and the drag from the air. The formula to find terminal velocity is:

\begin{equation}
V_t = \sqrt{\frac{4 d \rho_s g}{3 C_D \rho}}
\end{equation}

In this equation, \(g\) represents the acceleration due to gravity, \(C_D\) is the drag coefficient, and \(\rho\) is the air density. For a sphere in turbulent air flow, like a falling hailstone, \(C_D\) is usually about 0.47. Using the known values, we calculate:

\begin{equation}
V_t = \sqrt{\frac{4 \times 0.006 \times 917 \times 9.81}{3 \times 0.47 \times 1.2}}
\end{equation}
\begin{equation}
V_t = 11.29 \text{ m/s}
\end{equation}

\subsubsection*{Reynolds Number Calculation:}
The Reynolds number (\(Re\)) helps predict whether the flow around an object is smooth (laminar) or chaotic (turbulent). It compares the force of inertia to the force of viscosity in the fluid. The formula is:

\begin{equation}
Re = \frac{V_t \cdot d}{\nu}
\end{equation}

Here, \(\nu\) is the air's kinematic viscosity. For the conditions we're considering, \(\nu\) is about \(15 \times 10^{-6} \text{ m}^2/\text{s}\). Substituting the values, we get:

\begin{equation}
Re = \frac{11.29 \cdot 0.006}{15 \times 10^{-6}} = 4516
\end{equation}

This high \(Re\) indicates turbulent flow, as it is much larger than the critical range of 300 to 400 for spheres.

\subsubsection*{Calculating the Nusselt Number:}
The Nusselt number (\(Nu\)) measures how effectively heat is transferred from the hailstone to the air. It relates the convective to the conductive heat transfer. We use this equation:

\begin{equation}
Nu = 2 + (0.4 \cdot Re^{0.5} + 0.06 \cdot Pr^{0.67}) \cdot Pr^{0.4}
\end{equation}

The Prandtl number (\(Pr\)), which measures the relative thickness of the momentum and thermal layers in the fluid, for air is 0.713. Plugging in the values, we calculate:

\begin{equation}
Nu = 2 + (0.4 \cdot 4516^{0.5} + 0.06 \cdot 4516^{0.67}) \cdot 0.713^{0.4} = 40.2
\end{equation}

\subsubsection*{Heat Transfer Coefficient Calculation:}
Finally, the heat transfer coefficient (\(h\)) indicates the rate of heat transfer per unit area per unit temperature difference. It's connected to the Nusselt number by:

\begin{equation}
Nu = \frac{h \cdot d}{k}
\end{equation}

Using the previously calculated \(Nu\) and the thermal conductivity of air (\(k\)), we find \(h\) as:

\begin{equation}
h = \frac{Nu \cdot k}{d}
\end{equation}

\begin{equation}
h = \frac{40.2 \cdot 0.0257}{0.006} = 172.19 \text{ W/m}^2\text{K}
\end{equation}

This completes our calculations. The heat transfer coefficient for a 6 mm-diameter hailstone falling at terminal velocity in air at 287 K and 1 atm is \( 172.19 \) W/m\(^2\)K.

\newpage
\section{\#7}

Consider a counterflow, concentric tube heat exchanger that is designed to heat water from 20 to 80 \( ^{\circ}\)C. Hot oil is used as the heating agent, which enters the annulus at 160 \( ^{\circ}\)C and leaves at 140 \( ^{\circ}\)C. The inner tube, through which water flows, is thin-walled with a diameter of 20 mm. The overall heat transfer coefficient is given as 500 W/m\(^2\)K, and the design calls for a total heat transfer rate of 3000 W.

\subsection*{a. Determining the Length of the Heat Exchanger}

To determine the necessary length of the heat exchanger, we start by calculating the logarithmic mean temperature difference (LMTD), which adjusts for the variation in temperature across the heat exchanger. The formula for LMTD is:

\begin{equation}
\Delta T_{lm} = \frac{\Delta T_{o} - \Delta T_{i}}{\ln\left(\frac{\Delta T_{o}}{\Delta T_{i}}\right)}
\end{equation}

where:
\begin{itemize}
    \item \( \Delta T_{o} \) represents the temperature difference between the hot and cold fluids at the outlet.
    \item \( \Delta T_{i} \) represents the temperature difference at the inlet.
\end{itemize}

The temperatures are given as:
\begin{itemize}
    \item \( T_{h,i} = 160 \, ^{\circ}C \) is the inlet temperature of the hot oil.
    \item \( T_{h,o} = 140 \, ^{\circ}C \) is the outlet temperature of the hot oil.
    \item \( T_{c,o} = 80 \, ^{\circ}C \) is the outlet temperature of the water.
    \item \( T_{c,i} = 20 \, ^{\circ}C \) is the inlet temperature of the water.
\end{itemize}

Calculating the temperature differences:
\begin{equation}
\Delta T_{i} = T_{h,i} - T_{c,o} = 160 - 80 = 80 \, ^{\circ}C
\end{equation}
\begin{equation}
\Delta T_{o} = T_{h,o} - T_{c,i} = 140 - 20 = 120 \, ^{\circ}C
\end{equation}

By substituting these values into the LMTD equation, we obtain:
\begin{equation}
\Delta T_{lm} = \frac{120 - 80}{\ln(120/80)} = 98.65 \, ^{\circ}C
\end{equation}

The heat transfer rate equation \( q = U \times A \times \Delta T_{lm} \) allows us to solve for the area \( A \), and subsequently the length \( L \) since the area depends on length (\( A = \pi D_i L \)):
\begin{equation}
L = \frac{q}{U \times \pi \times D_{i} \times \Delta T_{lm}} = \frac{3000 \, \text{W}}{500 \, \text{W/m}^2\text{K} \times \pi \times 0.02 \, \text{m} \times 98.65 \, \text{K}} = 0.968 \, \text{m}
\end{equation}


\subsection*{b. Performance Degradation Due to Fouling}
\subsubsection*{Calculating the New Heat Transfer Rate (\(q_{new}\))}
To calculate the new heat transfer rate after fouling, use the formula:
\begin{equation}
q_{new} = \dot{m}_{c} C_{p,c} (T_{c,o} - T_{c,i}) = 50 \times (65 - 20) = 2250 \, \text{W}
\end{equation}

\subsubsection*{Calculating the New Outlet Temperature of the Oil (\(T_{h,o,new}\))}
The new outlet temperature of the oil can be calculated by adjusting for the new heat transfer rate:
\begin{equation}
T_{h,o,new} = T_{h,i} - \frac{q_{new}}{\dot{m}_{h} C_{p,h}} = 160 - \frac{2250 \, \text{W}}{150 \, \text{W/K}} = 145 \, ^{\circ}C
\end{equation}

\subsubsection*{Recalculating Temperature Differences and LMTD}
Recalculate the temperature differences at the inlet and outlet and the Logarithmic Mean Temperature Difference (LMTD) with these new temperatures:
\begin{align}
\Delta T_{i,new} &= T_{h,i} - T_{c,o,new} = 160 - 65 = 95 \, ^{\circ}C \\
\Delta T_{o,new} &= T_{h,o,new} - T_{c,i} = 145 - 20 = 125 \, ^{\circ}C \\
\Delta T_{lm,new} &= \frac{125 - 95}{\ln(125/95)} = 109.31 \, ^{\circ}C
\end{align}

\subsubsection*{Calculating the New Overall Heat Transfer Coefficient (\(U_{new}\))}
With the new LMTD, calculate the new overall heat transfer coefficient:
\begin{equation}
U_{new} = \frac{q_{new}}{A \times \Delta T_{lm,new}} = \frac{2250 \, \text{W}}{\pi \times 0.02 \, \text{m} \times 0.968 \, \text{m} \times 109.31 \, \text{K}} = 338.42 \, \text{W/m}^2\text{K}
\end{equation}

\subsubsection*{Calculating the Fouling Factor (\(R''_{f,c}\))}
Finally, find the water-side fouling factor by using the new overall heat transfer coefficient:
\begin{equation}
R''_{f,c} = \frac{1}{U_{new}} - \frac{1}{U} = \frac{1}{338.42} - \frac{1}{500} = 0.00095489 \, \text{m}^2\text{K/W}
\end{equation}

\end{document}